\documentclass[11pt]{article}
% 커널 SVM 퀴즈: 왜 w = Σ α_i x_i 이고, 왜 α 와 K 만으로 학습되나
% 빌드 (이 폴더에서):  tectonic svm_kernel_quiz_solution.tex   -> 정답 (빈칸에 파란 답)
%                      tectonic svm_kernel_quiz.tex            -> 학생용 (빈칸)
% 숫자는 덱과 같은 점 20개 예 (서포트 벡터 3개, α = 0.74, -0.37, -0.37, b = 0.04)

\usepackage[a4paper,margin=2.2cm]{geometry}
\usepackage{kotex}
\usepackage{amsmath,amssymb}
\usepackage{enumitem}
\usepackage{comment}
\usepackage{xcolor}
\usepackage{fancyhdr}
\usepackage{tikz}
\usepackage{graphicx}
\graphicspath{{../figs/}}
\usetikzlibrary{arrows.meta}

\ifdefined\STUDENT
  \excludecomment{answer}
  \newcommand{\blank}[2]{\underline{\hspace{#1}}}
  \newcommand{\pick}[2]{\text{(\ #1\ )}}
  \newcommand{\abox}[1]{\par\noindent\framebox[\dimexpr\linewidth-2\fboxsep-2\fboxrule\relax]{\rule{0pt}{#1}}\par}
\else
  \specialcomment{answer}{\begingroup\par\vspace{2pt}\color{blue}\small\setlength{\parskip}{2pt}\setlength{\parindent}{0pt}\noindent\textbf{풀이.}\ }{\par\endgroup}
  \newcommand{\blank}[2]{\underline{\makebox[#1]{\color{blue}#2}}}
  \newcommand{\pick}[2]{\text{(\ {\color{blue}\textbf{#2}}\ )}}
  \newcommand{\abox}[1]{}
\fi

\setlength{\headheight}{15pt}
\pagestyle{fancy}
\fancyhf{}
\fancyhead[L]{\small ML1 02-3 커널 SVM}
\fancyhead[R]{\bf 퀴즈: 왜 $\alpha$ 와 $K$ 만으로 되나\ifdefined\STUDENT\else\ \ (정답)\fi}
\fancyfoot[C]{\thepage}
\setlength{\parskip}{0.06in}
\setlength{\parindent}{0pt}

\newcounter{q}
\newcommand{\q}[1]{\par\vspace{10pt}\refstepcounter{q}{\large\bf Q\theq.}\ {\bf #1}\par\vspace{2pt}}
\newenvironment{parts}{\begin{enumerate}[label=(\alph*),leftmargin=2.2em,itemsep=6pt,topsep=2pt]}{\end{enumerate}}

\newcommand{\bw}{\mathbf{w}}
\newcommand{\bx}{\mathbf{x}}
\newcommand{\bu}{\mathbf{u}}
\newcommand{\bv}{\mathbf{v}}

\begin{document}

\noindent\fbox{\parbox{\dimexpr\linewidth-2\fboxsep-2\fboxrule}{
{\bf 약속} (덱과 같다): 경계 $\bw\cdot\bx + b = 0$, \ $\lVert\bw\rVert = 1$ 이면 점 $i$ 의 값 $y_i(\bw\cdot\bx_i + b)$ = 경계까지 거리 ($y_i = +1$ 빨강, $-1$ 파랑).\\
SVM: \ $\lVert\bw\rVert = 1$, 모든 $i$ 에서 $y_i(\bw\cdot\bx_i + b) \ge m$ 일 때 $m$ 을 최대로 하는 $\bw,\ b$.
}}

%%%%%%%%%%%%%%%%%%%%%%%%%%%%%%%%%%%%%%%%%%%%%%%%%%%%%%%%%%%%%%%%%%%%
\q{몸풀기: $\alpha$ 로 $\bw$ 만들기}
덱의 점 20개로 학습한 SVM (그림). 진한 점 3개가 서포트 벡터이고 나머지 17개는 $\alpha = 0$, \ $b = 0.04$.
\begin{center}\includegraphics[width=0.62\linewidth]{kn_quiz1.png}\end{center}
\vspace{-6pt}
\begin{parts}
\item $\sum_i \alpha_i \bx_i = \big(\,\blank{1.4cm}{0.70},\ \blank{1.4cm}{0.70}\,\big)$ \quad (소수 둘째 자리) \quad $\to$ 이것이 $\bw$ 다.
\item 새 점 $\bx = (0.5,\ -0.3)$: 서포트 벡터마다 $\alpha_i\, \bx\cdot\bx_i$ 를 구하면
\[ \blank{1.4cm}{0.21} \ + \ \blank{1.4cm}{0.07} \ + \ \blank{1.4cm}{-0.14}, \qquad f(\bx) = b + \textstyle\sum_i \alpha_i\, \bx\cdot\bx_i = \blank{1.4cm}{0.18} \ \to \ \pick{빨강 / 파랑}{빨강} \]
\item 한 번에: \ $\bw\cdot\bx + b = 0.70(0.5) + 0.70(-0.3) + 0.04 = \blank{1.4cm}{0.18}$ \quad (b) 와 \pick{같다 / 다르다}{같다}
\end{parts}
\begin{answer}
(a) $0.74(0.7) - 0.37(-0.6) - 0.37(0.1) = 0.518 + 0.222 - 0.037 = 0.70$, \ $0.74(0.2) - 0.37(-0.4) - 0.37(-1.1) = 0.148 + 0.148 + 0.407 = 0.70$.\\
(b) 점곱 $0.29,\ -0.18,\ 0.38$ 에 $\alpha$ 를 곱한다. $0.04 + 0.21 + 0.07 - 0.14 = 0.18 > 0$.
\end{answer}

%%%%%%%%%%%%%%%%%%%%%%%%%%%%%%%%%%%%%%%%%%%%%%%%%%%%%%%%%%%%%%%%%%%%
\clearpage
\q{증명: 최적의 $\bw$ 는 항상 학습 점의 합}
Q1 의 $\bw$ 가 학습 점의 합인 것은 우연이 아니다. $\bw$ 를 두 조각으로 나눈다.
\[ \bw = \bu + \bv, \qquad \bu = \textstyle\sum_i \alpha_i \bx_i \ \text{(학습 점으로 만든 부분)}, \qquad \bv\cdot\bx_i = 0 \ \text{(모든 학습 점과 수직)} \]
\begin{center}
\begin{tikzpicture}[scale=1.0, >={Stealth[length=2.2mm]}]
  \fill[blue!6] (0,0) -- (4.2,0) -- (5.4,1.4) -- (1.2,1.4) -- cycle;
  \draw[blue!40] (0,0) -- (4.2,0) -- (5.4,1.4) -- (1.2,1.4) -- cycle;
  \node[blue!60!black, font=\small] at (2.7,-0.3) {학습 점들이 만드는 평면};
  \foreach \p in {(0.9,0.4),(1.6,0.9),(2.3,0.3),(3.3,1.0),(3.9,0.5)} \fill[black!55] \p circle (1.6pt);
  \coordinate (O) at (1.9,0.55); \coordinate (U) at (3.4,0.95); \coordinate (W) at (3.4,2.35);
  \draw[->, very thick, blue!70!black] (O) -- (U) node[midway, below, font=\small] {$\bu$};
  \draw[->, very thick, red!70!black] (U) -- (W) node[midway, right, font=\small] {$\bv$};
  \draw[->, very thick] (O) -- (W) node[midway, left, font=\small] {$\bw$};
  \draw (3.4,1.15) -- (3.25,1.11) -- (3.25,0.91);
\end{tikzpicture}
\end{center}
\begin{parts}
\item $\bw\cdot\bx_i = \bu\cdot\bx_i + \bv\cdot\bx_i = \bu\cdot\bx_i + \blank{1.2cm}{0}$ \\
$\to$ 학습 점의 값 $y_i(\bw\cdot\bx_i + b)$ 는 \blank{1.2cm}{$\bu$} 와 $b$ 만으로 정해진다. $\bv$ 는 값을 \pick{바꾼다 / 못 바꾼다}{못 바꾼다}.
\item $\bu\cdot\bv = \sum_i \alpha_i\,(\bx_i\cdot\bv) = \blank{1.2cm}{0}$ \quad 이므로 \quad $\lVert\bw\rVert^2 = \lVert\bu\rVert^2 + \blank{1.2cm}{$\lVert\bv\rVert^2$}$ \quad (피타고라스)
\item $\lVert\bw\rVert = 1$ 인데 $\bv \ne \mathbf{0}$ 이면 \ $\lVert\bu\rVert$ \pick{$<$ / $=$ / $>$}{$<$} $1$.\\
$\bw,\ b$ 대신 $\bu/\lVert\bu\rVert,\ b/\lVert\bu\rVert$ 를 쓰면 길이는 1 이고, 모든 점의 값이 \blank{1.6cm}{$1/\lVert\bu\rVert$} 배\\ $\to$ 마진이 \pick{커진다 / 작아진다}{커진다}.\\
그래서 최적의 $\bw$ 에서는 $\bv = \blank{1.2cm}{$\mathbf{0}$}$, \ 즉 \ $\bw = \blank{2.2cm}{$\sum_i \alpha_i \bx_i$}$.
\item $\bw = \sum_i \alpha_i \bx_i$ 를 넣으면\\[2pt]
\hspace*{1em} 값 $= y_j\big(\,\blank{3cm}{$\sum_i \alpha_i\, \bx_i\cdot\bx_j$} + b\big)$, \qquad $\lVert\bw\rVert^2 = \blank{3.4cm}{$\sum_i\sum_j \alpha_i\alpha_j\, \bx_i\cdot\bx_j$}$\\[2pt]
$\to$ 학습에 데이터는 학습 점끼리의 \blank{1.4cm}{점곱} 으로만 들어간다.
\end{parts}
\begin{answer}
$\bv$ 는 어떤 점의 값에도 도움이 안 되면서 길이 1 중 일부를 잡아먹는다. 그러니 최적해에는 $\bv$ 가 없다. 소프트 마진도 길 안 점의 벌점이 $\bw\cdot\bx_i$ 로만 정해지므로 똑같다.
\end{answer}

%%%%%%%%%%%%%%%%%%%%%%%%%%%%%%%%%%%%%%%%%%%%%%%%%%%%%%%%%%%%%%%%%%%%
\clearpage
\q{커널로: 왜 $K$ 로 쓸 수 있나}
점 $\bx$ 를 특징 $\varphi(\bx)$ 로 옮긴 뒤, 그 위에서 선형 SVM: \ $f(\bx) = \bw\cdot\varphi(\bx) + b$. \ 빈칸을 채워라.
\begin{parts}
\item {\bf $\bw$ 는 학습 점의 합.} \ Q2 를 $\bx_i$ 대신 $\varphi(\bx_i)$ 로 반복하면 \ $\bw = \sum_i \alpha_i\, \blank{1.6cm}{$\varphi(\bx_i)$}$
\item {\bf 넣고 점곱을 나눈다.} (점곱은 더하기에 대해 나눠서 할 수 있다)
\[ f(\bx) = \Big(\textstyle\sum_i \alpha_i\, \varphi(\bx_i)\Big)\cdot\varphi(\bx) + b = \textstyle\sum_i \alpha_i\, \Big(\,\blank{2.6cm}{$\varphi(\bx_i)\cdot\varphi(\bx)$}\,\Big) + b \]
$\to$ $\varphi$ 는 늘 두 점을 옮긴 뒤의 \blank{1.4cm}{점곱} 모양으로만 나온다.
\item {\bf 그 점곱에 이름을 붙인다.} \ $K(\bx, \bx') = \varphi(\bx)\cdot\varphi(\bx')$ 이면
\[ f(\bx) = b + \textstyle\sum_i \alpha_i\, \blank{2cm}{$K(\bx_i, \bx)$} \]
\item {\bf 학습도 같다.} \ Q2 (d) 의 점곱 자리에 $K$ 를 넣으면
\[ \text{값} = y_j\Big(\textstyle\sum_i \alpha_i\, \blank{2cm}{$K(\bx_i, \bx_j)$} + b\Big), \qquad \lVert\bw\rVert^2 = \blank{3.4cm}{$\sum_i\sum_j \alpha_i\alpha_j K(\bx_i, \bx_j)$} \]
$\to$ 학습과 분류 어디에도 $\varphi$ 를 \pick{만들어야 한다 / 안 만들어도 된다}{안 만들어도 된다}.
\item RBF 는 $\varphi$ 의 특징이 \blank{1.2cm}{$\infty$} 개라 직접 만들 수 없다. 그래도 학습 점이 $n$ 개면 숫자 $n$ 개 ($\alpha$) 로 충분한 이유를 (a) 를 이용해 한 문장으로 써라.
\abox{1.6cm}
\end{parts}
\begin{answer}
한 줄 요약: $\bw$ 가 점들의 합 $\to$ $\bw\cdot\varphi(\bx)$ 가 점곱들의 합 $\to$ 점곱 자리에 $K$.\\
(e) 최적의 $\bw$ 는 $n$ 개 점 $\varphi(\bx_1), \dots, \varphi(\bx_n)$ 의 합이다 (Q2 에서 $\bv = \mathbf{0}$). 특징이 무한 개여도 $\alpha_1, \dots, \alpha_n$ 만 알면 된다.
\end{answer}

%%%%%%%%%%%%%%%%%%%%%%%%%%%%%%%%%%%%%%%%%%%%%%%%%%%%%%%%%%%%%%%%%%%%
\clearpage
\q{확인: 다항 커널 $(1 + \bx\cdot\bx')^2$ 은 정말 $\varphi$ 의 점곱인가}
2차원 점 $\bx = (x_1, x_2)$ 를 특징 6개로 옮긴다: \quad $\varphi(\bx) = \big(1,\ \sqrt2 x_1,\ \sqrt2 x_2,\ x_1^2,\ \sqrt2 x_1x_2,\ x_2^2\big)$.
\begin{parts}
\item 6개를 짝지어 곱해 더하면
\[ \varphi(\bx)\cdot\varphi(\bx') = 1 + 2x_1x_1' + 2x_2x_2' + x_1^2x_1'^2 + \blank{1.8cm}{$2x_1x_2x_1'x_2'$} + x_2^2x_2'^2 \]
\item 묶으면 \quad $= 1 + 2\,(x_1x_1' + x_2x_2') + (x_1x_1' + x_2x_2')^2 = 1 + 2\,(\,\blank{1.4cm}{$\bx\cdot\bx'$}\,) + (\,\blank{1.4cm}{$\bx\cdot\bx'$}\,)^2 = \blank{2.2cm}{$(1 + \bx\cdot\bx')^2$}$
\item 숫자로: \ $\bx = (1, 2)$, \ $\bx' = (3, 1)$.
\begin{itemize}[leftmargin=1.5em,itemsep=3pt]
\item 커널로 (2차원 점곱 한 번): \ $\bx\cdot\bx' = \blank{1cm}{5}$, \quad $K = (1 + \bx\cdot\bx')^2 = \blank{1cm}{36}$
\item 옮겨서 (6차원 점곱): \ $\varphi(1,2) = (1,\ \sqrt2,\ 2\sqrt2,\ 1,\ 2\sqrt2,\ 4)$,\\[2pt] \hspace*{7.6em}$\varphi(3,1) = \big(\,\blank{4.2cm}{$1,\ 3\sqrt2,\ \sqrt2,\ 9,\ 3\sqrt2,\ 1$}\,\big)$\\[2pt]
$\varphi(1,2)\cdot\varphi(3,1) = 1 + 6 + \blank{0.8cm}{4} + \blank{0.8cm}{9} + \blank{0.8cm}{12} + 4 = \blank{1cm}{36}$
\end{itemize}
\item 두 값은 \pick{같다 / 다르다}{같다}. 커널은 특징 6개를 \pick{만든다 / 안 만든다}{안 만든다}.
\end{parts}
\begin{answer}
$\sqrt2 x_1 \cdot \sqrt2 x_1' = 2x_1x_1'$ 처럼 $\sqrt2$ 끼리 곱해 2 가 된다. 그래서 $\sqrt2$ 가 붙어 있어야 $(1 + \bx\cdot\bx')^2$ 과 정확히 같다.\\
$2\sqrt2 \cdot \sqrt2 = 4$, \ $1 \cdot 9 = 9$, \ $2\sqrt2 \cdot 3\sqrt2 = 12$.
\end{answer}

\end{document}
